\documentclass[10pt]{article}
\usepackage[paperwidth=190mm,paperheight=64mm,margin=1mm]{geometry}
\usepackage{fontspec,booktabs,multirow,graphicx,amsmath}
\usepackage[table]{xcolor}
\IfFontExistsTF{Times New Roman}{\setmainfont{Times New Roman}}{\setmainfont{TeX Gyre Termes}}\pagestyle{empty}\setlength{\parindent}{0pt}
\definecolor{grayband}{gray}{.9}\definecolor{lavender}{RGB}{244,241,252}\definecolor{greenband}{RGB}{240,255,240}
\setlength{\tabcolsep}{3.2pt}\renewcommand{\arraystretch}{1.0}
\begin{document}
\graphicspath{{assets/}}
{\fontsize{14}{15.5}\selectfont Table 8: \textbf{Representation generation ablation experiments.} The default RDM backbone is of 12 blocks and 1536 hidden dimensions, trained for 100 epochs, and takes 250 sampling steps during generation. The representation Fr\'echet Distance (rep FD) is evaluated between 50K generated representations and representations extracted from the ImageNet training set by MoCo v3 ViT-B. Default settings are marked in \colorbox{grayband}{gray}.\par}\vspace{1.5mm}
\noindent\makebox[\linewidth][c]{%
\begin{minipage}[t]{.24\linewidth}\centering\fontsize{8}{9.1}\selectfont\begin{tabular}{cccc}
\#Blocks & FID & IS & rep FD \\
\midrule
3 & 7.53 & 113.5 & 0.71 \\
6 & 5.40 & 132.9 & 0.53 \\
12 & \cellcolor{grayband}\textbf{5.07} & \cellcolor{grayband}\textbf{142.5} & \cellcolor{grayband}\textbf{0.48} \\
18 & 5.20 & 141.9 & 0.50 \\
24 & 5.13 & 141.5 & 0.49 \\
\end{tabular}\par\vspace{1mm}{\fontsize{9}{10.3}\selectfont\raggedright (a) \textbf{Model depth.} A deeper RDM can improve generation performance.\par}\end{minipage}\hspace{1.5mm}%
\begin{minipage}[t]{.24\linewidth}\centering\fontsize{8}{9.1}\selectfont\begin{tabular}{cccc}
Hidden Dim & FID & IS & rep FD \\
\midrule
256 & 12.99 & 67.3 & 5.98 \\
512 & 9.07 & 99.8 & 1.19 \\
1024 & 5.35 & 132.0 & 0.56 \\
1536 & \cellcolor{grayband}\textbf{5.07} & \cellcolor{grayband}142.5 & \cellcolor{grayband}\textbf{0.48} \\
2048 & 5.09 & \textbf{142.8} & \textbf{0.48} \\
\end{tabular}\par\vspace{1mm}{\fontsize{9}{10.3}\selectfont\raggedright (b) \textbf{Model width.} A wider RDM can improve generation performance.\par}\end{minipage}\hspace{1.5mm}%
\begin{minipage}[t]{.24\linewidth}\centering\fontsize{8}{9.1}\selectfont\begin{tabular}{cccc}
Epochs & FID & IS & rep FD \\
\midrule
10 & 5.94 & 124.4 & 0.87 \\
50 & 5.21 & 138.3 & 0.54 \\
100 & \cellcolor{grayband}5.07 & \cellcolor{grayband}142.5 & \cellcolor{grayband}0.48 \\
200 & 5.07 & \textbf{145.1} & \textbf{0.47} \\
300 & \textbf{5.05} & 144.3 & \textbf{0.47} \\
\end{tabular}\par\vspace{1mm}{\fontsize{9}{10.3}\selectfont\raggedright (c) \textbf{Training epochs.} Training RDM longer improves generation performance.\par}\end{minipage}\hspace{1.5mm}%
\begin{minipage}[t]{.24\linewidth}\centering\fontsize{8}{9.1}\selectfont\begin{tabular}{cccc}
\#Steps & FID & IS & rep FD \\
\midrule
20 & 5.80 & 120.3 & 0.87 \\
50 & 5.28 & 133.0 & 0.55 \\
100 & 5.15 & 138.1 & \textbf{0.48} \\
250 & \cellcolor{grayband}\textbf{5.07} & \cellcolor{grayband}\textbf{142.5} & \cellcolor{grayband}\textbf{0.48} \\
500 & \textbf{5.07} & \textbf{142.9} & 0.49 \\
\end{tabular}\par\vspace{1mm}{\fontsize{9}{10.3}\selectfont\raggedright (d) \textbf{Diffusion steps.} More sampling steps can improve generation performance.\par}\end{minipage}}

\end{document}
